\section{Introduction}
Android is the dominant mobile operating system and acts as a primary gateway to the digital world for billions of users. It is crucial for the core operating system to be a reliable, stable, and trustworthy platform that allows users to run applications and services of their choice. While third-party application privacy has been widely studied \cite{appDataCollection1,appDataCollection2,appDataCollection3,appDataCollection4}, less attention is typically paid to the operating system layer and to privileged service frameworks that silently enable additional functionality.

While the Android core is the Android Open Source Project (AOSP), most consumer devices include GMS (Google Mobile Services) layered on top. They are implemented as privileged system applications to provide baseline functionality. Application developers commonly assume their presence by default. This design tightly couples many functions into a single service stack. Enabling useful capabilities may implicitly enable others that a user would not choose, such as advertising identifiers, telemetry, and behavioral tracking. Because Google’s business model is strongly advertising-driven, many users may be unwilling to entrust this provider with broad access to their devices and data.

In principle, users should be able to choose which services they rely on and have confidence that these choices are enforced by the platform’s permission and isolation mechanisms. This demand for transparency and controllability motivates alternative approaches such as microG\footnote{\url{https://github.com/microg/GmsCore/wiki}} and GrapheneOS sandboxed Google Play\footnote{\url{https://grapheneos.org/usage\#sandboxed-google-play}}. These solutions aim to preserve device functionality and application compatibility while reducing privilege and improving user control over Google-related components.

Privacy-focused Android operating systems aim to reduce data exposure at the OS level. They achieve this by limiting preinstalled privileged components, tightening default settings, and providing more transparent controls. We focus on three representative approaches: GrapheneOS\footnote{\url{https://grapheneos.org/}} prioritizes a hardened security model and strong isolation. Google components are not included by default, but can be installed in a sandboxed mode as ordinary applications. LineageOS\footnote{\url{https://lineageos.org/}} provides a lightweight, broadly supported AOSP-based system with minimal preinstalled software. It can be deployed in three configurations: without Google services, with Gapps\footnote{\url{https://mindthegapps.com/}} (proprietary Google Play Services), or with microG\footnote{\url{https://lineage.microg.org/}}. Finally, iod\'eOS\footnote{\url{https://iode.tech/iodeos/}} is a privacy-oriented distribution that combines a de-Googled baseline with additional privacy features and preconfigured app ecosystem choices, shipping with microG.

Prior work has established that Android privacy risks arise not only from user-installed applications, but also from the operating system and its preinstalled components \cite{android-os-snooping2021,android-os-snooping2023,google-dialer-messages,cookies-identifiers2025,google-data-collection2018}. Official Android builds transmit substantial data to OS vendors and third parties, including telemetry, location, app activity, and persistent identifiers. Google serves as the dominant driver of this background tracking ecosystem. Prior literature shows that removing or replacing these components with open-source variants drastically reduces unsolicited communication \cite{android-os-snooping2021}. However, alternative distributions remain underexplored in reproducible measurement studies. This motivates our systematic comparison of stock and privacy-oriented systems.

Our main contributions are as follows: first, we present a comparative network-level privacy analysis of six Android configurations across diverse Google integration models; second, we provide the first systematic empirical comparison between privileged GMS, sandboxed Google Play Services, and microG under identical conditions; third, we design a controlled, real-device traffic interception and TLS decryption measurement methodology tailored for Android 16; and fourth, we quantify how alternative distributions reduce background data and identifier exposure while preserving core functionality.
